\documentclass[11pt,a4paper]{article}

% --- PACKAGES ---
\usepackage[utf8]{inputenc}
\usepackage[indonesian]{babel}
\usepackage{amsmath, amssymb, amsfonts}
\usepackage{geometry}
\geometry{margin=1in,headheight=14pt}
\usepackage{graphicx}
\usepackage{booktabs}
\usepackage{hyperref}
\usepackage{tcolorbox}
\tcbuselibrary{skins,breakable}
\usepackage{listings}
\usepackage{xcolor}
\usepackage{fancyhdr}

% --- STYLING & COLORS ---
\definecolor{unibblue}{RGB}{0, 51, 102}
\definecolor{unibgold}{RGB}{212, 175, 55}
\definecolor{darkgreen}{RGB}{34, 139, 34}

\hypersetup{
    colorlinks=true,
    linkcolor=unibblue,
    citecolor=darkgreen,
    urlcolor=blue
}

\pagestyle{fancy}
\fancyhf{}
\rhead{\textbf{TKE-41XX: Optimisasi untuk Teknik Elektro}}
\lhead{Catatan Kuliah Minggu ke-13}
\rfoot{Halaman \thepage}

\lstset{
  language=Python,
  basicstyle=\ttfamily\footnotesize,
  keywordstyle=\color{unibblue}\bfseries,
  commentstyle=\color{darkgreen}\itshape,
  stringstyle=\color{red!70!black},
  numbers=left,
  numberstyle=\tiny\color{gray},
  numbersep=5pt,
  breaklines=true,
  frame=single,
  rulecolor=\color{unibblue!30},
  tabsize=4
}

\title{\textbf{Catatan Kuliah Minggu ke-13:}\\Economic Dispatch Sistem Tenaga Listrik Berbasis Metaheuristik Particle Swarm Optimization (PSO)}
\author{\textbf{Ir. Novalio Daratha S.T., M.Sc., Ph.D.}\\Program Studi S1 Teknik Elektro --- Universitas Bengkulu}
\date{Semester Ganjil 2026}

\begin{document}

\maketitle

\begin{tcolorbox}[colback=unibblue!5,colframe=unibblue,title=\textbf{Ringkasan Eksekutif dan CPMK}]
Dokumen ini menyajikan panduan perkuliahan komprehensif mengenai aplikasi \textit{Particle Swarm Optimization} (PSO) pada masalah \textit{Economic Dispatch} (ED) dalam sistem tenaga listrik. Mahasiswa mempelajari pembentukan fungsi biaya bahan bakar pembangkit termal (termasuk efek \textit{valve-point loading}), perumusan kendala keseimbangan daya dan batas generator, serta mekanisme pengkodean partikel PSO dan *constraint handling*.
\end{tcolorbox}

\tableofcontents
\newpage

\section{Pendahuluan Economic Dispatch}
Dalam pengoperasian sistem tenaga listrik modern, penyediaan energi listrik kepada konsumen harus memenuhi kriteria utama: **andal, berkualitas, dan ekonomis**. Masalah \textit{Economic Dispatch} (ED) berada pada tingkatan perencanaan operasi jangka pendek (*short-term operational planning*) yang bertujuan menentukan alokasi daya aktif ($P_{G,i}$) untuk setiap unit pembangkit termal yang sedang beroperasi sedemikian rupa sehingga total biaya bahan bakar (\$/jam) mencapai nilai minimum.

\subsection{Karakteristik Biaya Bahan Bakar Pembangkit Termal}
Fungsi biaya bahan bakar pembangkit termal pada kondisi ideal umumnya didekati dengan persamaan kuadratik:
\begin{equation}
    F_i(P_{G,i}) = a_i P_{G,i}^2 + b_i P_{G,i} + c_i \quad [\text{\$/jam}]
\end{equation}
di mana:
\begin{itemize}
    \item $P_{G,i}$: Daya aktif keluaran pembangkit ke-$i$ dalam MegaWatt (MW).
    \item $a_i$: Koefisien biaya bahan bakar kuadratik (\$/MW$^2$h).
    \item $b_i$: Koefisien biaya bahan bakar linier (\$/MWh).
    \item $c_i$: Biaya tanpa beban / *no-load cost* (\$/jam).
\end{itemize}

\subsection{Efek Valve-Point Loading (Non-Konveksitas)}
Pada kenyataannya, turbin uap menggunakan sekumpulan katup (*valves*) untuk mengatur masukan uap. Saat katup pembuka baru diaktifkan, terjadi efek pembukaan tercekik (*wire-drawing effect*) yang menimbulkan kerugian efisiensi sementara. Fenomena ini dimodelkan dengan menambahkan komponen sinusoidal pada fungsi biaya kuadratik:
\begin{equation}
    F_i(P_{G,i}) = a_i P_{G,i}^2 + b_i P_{G,i} + c_i + |e_i \sin(f_i(P_{i,\min} - P_{G,i}))|
\end{equation}
di mana $e_i$ dan $f_i$ adalah koefisien efek *valve-point*. Penambahan suku rippel ini membuat fungsi biaya bersifat **non-konveks, non-diferensiabel, dan memiliki banyak optimum lokal**, sehingga metode optimisasi berbasis turunan klasik tidak lagi efektif.

\section{Formulasi Matematis Masalah ED}

\subsection{Fungsi Tujuan}
Meminimalkan total biaya bahan bakar dari seluruh $N_g$ pembangkit:
\begin{equation}
    \min_{\mathbf{P}_G} F_T = \sum_{i=1}^{N_g} F_i(P_{G,i})
\end{equation}

\subsection{Kendala-Kendala (Constraints)}

\subsubsection{1. Keseimbangan Daya (Power Balance Constraint)}
Total daya pembangkitan harus sama dengan jumlah total beban sistem ($P_D$) dan rugi-rugi transmisi ($P_L$):
\begin{equation}
    \sum_{i=1}^{N_g} P_{G,i} - P_D - P_L = 0
\end{equation}

Rugi-rugi transmisi $P_L$ dimodelkan menggunakan matriks rugi-rugi Kron ($B$-coefficients):
\begin{equation}
    P_L = \sum_{i=1}^{N_g} \sum_{j=1}^{N_g} P_{G,i} B_{ij} P_{G,j} + \sum_{i=1}^{N_g} B_{0i} P_{G,i} + B_{00}
\end{equation}

\subsubsection{2. Batas Kapasitas Pembangkit (Generator Limits)}
Daya keluaran setiap unit pembangkatan dibatasi oleh kemampuan mekanis dan termal unit:
\begin{equation}
    P_{i,\min} \le P_{G,i} \le P_{i,\max}, \quad i = 1, 2, \dots, N_g
\end{equation}

\section{Penyelesaian dengan Particle Swarm Optimization (PSO)}

\subsection{Representasi Partikel}
Pada masalah ED dengan $N_g$ unit pembangkit, setiap partikel $k$ mewakili satu calon solusi alokasi daya:
\begin{equation}
    \mathbf{x}_k = [P_{G,1}^{(k)}, P_{G,2}^{(k)}, \dots, P_{G,N_g}^{(k)}]
\end{equation}
Kecepatan partikel mewakili laju perubahan alokasi daya:
\begin{equation}
    \mathbf{v}_k = [v_1^{(k)}, v_2^{(k)}, \dots, v_{N_g}^{(k)}]
\end{equation}

\subsection{Fungsi Fitness dan Penalti}
Untuk menangani kendala kesetaraan daya, digunakan fungsi penalti kuadratik:
\begin{equation}
    \text{Fitness}(\mathbf{x}_k) = \sum_{i=1}^{N_g} F_i(P_{G,i}^{(k)}) + \lambda_{\text{pen}} \left( \sum_{i=1}^{N_g} P_{G,i}^{(k)} - P_D - P_L \right)^2
\end{equation}
di mana $\lambda_{\text{pen}}$ adalah faktor penalti berbobot besar (misal $10^5$).

\subsection{Mekanisme Repair Generator (Slack Generator)}
Pendekatan alternatif yang sangat efisien untuk menangani kesetaraan daya adalah penetapan generator terakhir $N_g$ sebagai *slack generator*:
\begin{equation}
    P_{G,N_g} = P_D + P_L - \sum_{i=1}^{N_g-1} P_{G,i}
\end{equation}
Setelah nilai $P_{G,N_g}$ dihitung, dilakukan *clamping* batas daya:
\begin{equation}
    P_{G,N_g} = \min \left( P_{N_g,\max}, \max \left( P_{N_g,\min}, P_{G,N_g} \right) \right)
\end{equation}

\section{Implementasi Program Python}

Berikut adalah skrip Python lengkap untuk menyelesaikan Economic Dispatch 3 Pembangkit menggunakan PSO:

\begin{lstlisting}[caption={Skrip Python Economic Dispatch dengan PSO}]
import numpy as np

# Data Pembangkit: [a ($/MW^2h), b ($/MWh), c ($/h), Pmin (MW), Pmax (MW)]
gen_data = np.array([
    [0.00156, 7.92, 561, 100, 600],
    [0.00194, 7.85, 310, 100, 400],
    [0.00482, 7.97, 78, 50, 200]
])

P_D = 850.0  # Total beban (MW)
N_g = len(gen_data)

# Parameter PSO
N_particles = 40
Max_iter = 100
w = 0.7
c1 = 1.5
c2 = 1.5

# Inisialisasi partikel
Pmin = gen_data[:, 3]
Pmax = gen_data[:, 4]
X = np.random.uniform(Pmin, Pmax, size=(N_particles, N_g))
V = np.zeros((N_particles, N_g))

def evaluate_cost(x_vec):
    P = np.copy(x_vec)
    # Terapkan clamping pada N_g - 1 generator
    P[:-1] = np.clip(P[:-1], Pmin[:-1], Pmax[:-1])
    # Hitung generator slack
    P[-1] = P_D - np.sum(P[:-1])
    
    # Hitung penalti jika slack melanggar batas
    penalty = 0
    if P[-1] < Pmin[-1] or P[-1] > Pmax[-1]:
        penalty = 1e5 * (abs(P[-1] - np.clip(P[-1], Pmin[-1], Pmax[-1])))**2
        
    cost = np.sum(gen_data[:, 0] * P**2 + gen_data[:, 1] * P + gen_data[:, 2])
    return cost + penalty, P

# Evaluasi awal
pbest_X = np.copy(X)
pbest_fit = np.array([evaluate_cost(X[i])[0] for i in range(N_particles)])
gbest_idx = np.argmin(pbest_fit)
gbest_X = np.copy(pbest_X[gbest_idx])
gbest_fit = pbest_fit[gbest_idx]

# Loop Utama PSO
for t in range(Max_iter):
    for i in range(N_particles):
        r1, r2 = np.random.rand(N_g), np.random.rand(N_g)
        V[i] = w * V[i] + c1 * r1 * (pbest_X[i] - X[i]) + c2 * r2 * (gbest_X - X[i])
        X[i] = X[i] + V[i]
        
        fit_val, _ = evaluate_cost(X[i])
        if fit_val < pbest_fit[i]:
            pbest_fit[i] = fit_val
            pbest_X[i] = np.copy(X[i])
            if fit_val < gbest_fit:
                gbest_fit = fit_val
                gbest_X = np.copy(X[i])

_, P_opt = evaluate_cost(gbest_X)
print(f"Alokasi Daya Optimal: P1 = {P_opt[0]:.2f} MW, P2 = {P_opt[1]:.2f} MW, P3 = {P_opt[2]:.2f} MW")
print(f"Total Biaya Operasi Minimum: ${gbest_fit:.2f} / jam")
\end{lstlisting}

\section{Kesimpulan dan Diskusi}
1. Masalah \textit{Economic Dispatch} dapat diformulasikan secara matematis dan diselesaikan secara cepat menggunakan PSO.
2. Penggunaan strategi *slack generator repair* mempercepat konvergensi PSO dibandingkan penalti biasa.
3. Pendekatan berbasis PSO sangat adaptif untuk memperhitungkan non-konveksitas *valve-point loading* dan rugi-rugi transmisi $B_{ij}$.

\end{document}
